\lysh{20854}{4}
\begin{alist}
\item Η συνάρτηση $f$ είναι άρτια, διότι: \\
Για κάθε $x \in \mathbb{R}$, το $-x \in \mathbb{R}$ και ισχύει ότι $f(-x) = e^{|-x|} = e^{|x|} = f(x)$.
\item Αρκεί να αποδείξουμε ότι $f(x) \ge f(0)$ για κάθε $x \in \mathbb{R}$. Ισοδύναμα έχουμε ότι
\[e^{|x|} \ge e^0 \overset{e^x \uparrow}{\Longleftrightarrow} |x| \ge 0,\]
η οποία ισχύει για κάθε $x \in \mathbb{R}$. Η ισότητα ισχύει μόνο όταν $x = 0$. \\
Επομένως, η $f$ παρουσιάζει ελάχιστο στο $x=0$ και η ελάχιστη τιμή της είναι η $f(0) = 1$.
\item Αν $x \ge 0$, τότε $|x| = x$ και $f(x) = e^x$. \\
Αν $x < 0$, τότε $|x| = -x$ και $f(x) = e^{-x}$. \\
Έτσι προκύπτει η δίκλαδη συνάρτηση $f(x) = \begin{cases} e^{-x}, & x < 0 \\ e^x, & x \ge 0 \end{cases}$, η οποία σύμφωνα με το ερώτημα $\alpha)$ είναι άρτια. Οπότε είναι συμμετρική ως προς τον άξονα $y'y$. \\
Επομένως, αποτελείται από την γραφική παράσταση της $h(x) = e^x$, $x \ge 0$ και την συμμετρική της $h$ ως προς τον άξονα $y'y$ για $x < 0$. \\
Επιπλέον, από το ερώτημα $\beta)$ γνωρίζουμε ότι η $f$ παρουσιάζει ελάχιστο στο $x=0$ και η ελάχιστη τιμή της είναι η $f(0) = 1$. \\
Επομένως η γραφική παράσταση της $f$ δίνεται από το παρακάτω σχήμα:
\begin{center}
\begin{tikzpicture}
\begin{axis}[
    axis lines=middle,
    xmin=-3.5, xmax=3.5,
    ymin=-0.5, ymax=6.5,
    xtick={-3,-2,-1,0,1,2,3},
    ytick={1,2,3,4,5,6},
    grid=both,
    grid style={line width=.1pt, draw=gray!30},
    major grid style={line width=.2pt,draw=gray!50},
    thick,
    width=9cm, height=8cm,
    samples=100
]
\addplot[domain=0:1.85, black!70, ultra thick] {exp(x)} node[pos=0.8, right, black] {\large $C_f$};
\addplot[domain=-1.85:0, maincolor, ultra thick] {exp(-x)};
\addplot[domain=-3.5:0, dashed, ultra thick,secondarycolor ] {exp(x)};
\fill[black] (axis cs:0, 1) circle (2.5pt);
\end{axis}
\end{tikzpicture}
\end{center}
\end{alist}

